% !PRINTED 314
\weberpage{314}

Thus, if we set $e^{-\frac{5\pi i}{12}} = \varepsilon$ for brevity, we have the
following permutations belonging together:

\weq{21}{\begin{array}{cccccccc}
\omega, & u, & v_\infty, & v_0, & v_1, & v_2, & v_3, & v_4,\\
\omega+2, & e^{-\frac{\pi i}{12}}u, & \varepsilon v_\infty, & \varepsilon v_2, & \varepsilon v_3, & \varepsilon v_4, & \varepsilon v_0, & \varepsilon v_1,\\
-\dfrac{1}{\omega}, & u, & v_0, & v_\infty, & v_4, & v_2, & v_3, & v_1,\\
\dfrac{\omega-1}{\omega+1}, & \dfrac{\sqrt{2}}{u}, & -\dfrac{\sqrt{2}}{v_1}, & -\dfrac{\sqrt{2}}{v_4}, & -\dfrac{\sqrt{2}}{v_0}, & -\dfrac{\sqrt{2}}{v_2}, & -\dfrac{\sqrt{2}}{v_3}, & -\dfrac{\sqrt{2}}{v_\infty}.
\end{array}}

We now introduce the five quantities $w_z$ by the equation:

\weq{22}{w_z = \frac{(v_\infty - v_z)(v_{z+1}-v_{z-1})(v_{z+2}-v_{z-2})}{\sqrt{5}\,u^3},}

that is:

\[
\begin{aligned}
w_0 &= \frac{(v_\infty - v_0)(v_1-v_4)(v_2-v_3)}{\sqrt{5}\,u^3},\\
w_1 &= \frac{(v_\infty - v_1)(v_2-v_0)(v_3-v_4)}{\sqrt{5}\,u^3},\\
w_2 &= \frac{(v_\infty - v_2)(v_3-v_1)(v_4-v_0)}{\sqrt{5}\,u^3},\\
w_3 &= \frac{(v_\infty - v_3)(v_4-v_2)(v_0-v_1)}{\sqrt{5}\,u^3},\\
w_4 &= \frac{(v_\infty - v_4)(v_0-v_3)(v_1-v_2)}{\sqrt{5}\,u^3},
\end{aligned}
\]

so that, if one makes use in the last row of the fact that by (18) the product
of the six quantities $v_z$ has the value $u^6$, there follow from (21) these
permutations of the $w$ belonging together:

\weq{23}{\begin{array}{cccccc}
\omega, & w_0, & w_1, & w_2, & w_3, & w_4,\\
\omega+2, & -w_2, & -w_3, & -w_4, & -w_0, & -w_1,\\
-\dfrac{1}{\omega}, & w_0, & w_2, & w_1, & w_4, & w_3,\\
\dfrac{\omega-1}{\omega+1}, & -w_0, & -w_3, & -w_4, & -w_2, & -w_1.
\end{array}}

The functions $w$ cannot become infinite for any value of $u$ different from 0
and $\infty$. Let us therefore assume the equation whose roots are the five
quantities $w_z$ in the form:

\[
w^5 + A_1w^4 + A_2w^3 + A_3w^2 + A_4w + A_5 = 0,
\]
%CONT>>
